\section{Conclusion}

This paper presented an LLVM-based floating-point debugging framework that unifies Error-Free Transformations (EFTs) and condition-number analysis within a single shadow execution. Rather than maintaining separate instrumentation pipelines, our approach propagates a unified shadow state that simultaneously tracks accumulated numerical error through EFT/MPFR and numerical sensitivity through condition numbers. This shared representation enables the two analyses to complement one another while avoiding redundant instrumentation.

Our evaluation on the Herbie and PolyBench benchmark suites demonstrates that the combined analysis provides more informative diagnostics than EFT alone. In particular, incorporating condition numbers enables the framework to distinguish cancellation and sensitivity cases that cannot be reliably identified from propagated residuals alone, while preserving the efficient runtime characteristics of EFT-based shadow execution. On the Herbie benchmarks, the combined analysis agrees with Herbie's reported numerical issues across all evaluated benchmarks and achieves a median runtime overhead of approximately $2.5\times$ under \texttt{-O2}.

The current implementation also exposes several opportunities for future work. Extending the framework with dedicated overflow and underflow oracles, logarithmic shadow representations to distinguish sensitivity beyond the IEEE-754 dynamic range, and further reducing instrumentation overhead through compiler optimizations would improve both the accuracy and scalability of dynamic floating-point debugging.